% -----------------------------------------------------------------------------
% new_full_diagram.tex  --  one-pager sketch of the proximal prior pipeline.
%
% Deliverable T1.1 (fig:prior_pipeline) of docs/prior_appendix_plan.md.
%
% VISUAL GRAMMAR (kept consistent so the figure reads structurally):
%   * SQUARE node      = a parameter  (sampled dimension + fixed MLP: [d_X,f_X] ...)
%   * CIRCLE node      = a data / random variable  (X*,U*,W*,Z*, U^aug, A, (Y0,Y1))
%   * ROUNDED BOX      = a process step  (Generate, post-processing, Sample A/Y,
%                        aggregate, accept/reject, output)
%   * dotted arrow     = a parameter configures a step   ([d_X,f_X] ..> Generate X)
%   * solid arrow      = data flow (black generic; red into A; blue into Y)
%   * dashed teal      = confounding duplication  (U* -> U^aug -> {A,Y})
%   * orange dashed    = resample loop
%
% LAYOUT (top -> bottom): header -> parameter row -> generation row ->
%   post-processing box (star circles + caption) -> {Sample A | f_aug | Sample Y}
%   -> aggregate -> accept/reject -> output.
%
% X lives on T rows and is only *broadcast* across K, so the green T*K plate is
% L-SHAPED: its top-left notch excludes the Generate-X box and X*; U,W,Z,A,Y and
% the K-aggregate sit inside it.
%
% Code of record: proximalfm/prior/_dataset.py::SCMPrior.{get_batch,generate_dataset},
%                 proximalfm/prior/_reg2cls.py (process_features / dominates),
%                 proximalfm/prior/_prior_config.py (flag defaults).
%
% Bare tikzpicture fragment. Requires (all in docs/latex/preamble.tex):
%   tikz + libraries positioning, arrows.meta, calc.
% If it runs wider/taller than the text block, wrap the \input in
%   \resizebox{\textwidth}{!}{...}  or place it on a dedicated figure page.
%
% NOTE: custom styles are prefixed "pp" to avoid clashing with built-in TikZ keys.
% -----------------------------------------------------------------------------
\begin{tikzpicture}[
  line cap=round, line join=round,
  font=\footnotesize,
  % ---- process boxes (rounded rectangles) -----------------------------------
  pphdr/.style={rounded corners=2.5pt, draw=black!45, fill=black!3,
    line width=0.5pt, inner sep=5pt, align=center},
  ppgen/.style={rounded corners=2.5pt, draw=black!55, fill=gray!4,
    line width=0.5pt, inner sep=4pt, align=center},
  ppabox/.style={rounded corners=2.5pt, draw=red!60!black, fill=red!4,
    line width=0.6pt, inner sep=5pt, align=left, font=\scriptsize},
  ppybox/.style={rounded corners=2.5pt, draw=blue!60!black, fill=blue!5,
    line width=0.6pt, inner sep=5pt, align=left, font=\scriptsize},
  ppagg/.style={rounded corners=2.5pt, draw=brown!65!black, fill=brown!8,
    line width=0.6pt, inner sep=5pt, align=center, font=\scriptsize},
  ppacc/.style={rounded corners=2.5pt, draw=orange!72!black, fill=orange!6,
    line width=0.6pt, inner sep=5pt, align=center, font=\scriptsize},
  ppout/.style={rounded corners=2.5pt, draw=black!55, fill=gray!4,
    line width=0.6pt, inner sep=5pt, align=center},
  % ---- parameters (squares) -------------------------------------------------
  ppparam/.style={rectangle, draw=black!55, fill=yellow!12, line width=0.5pt,
    inner sep=3.5pt, font=\scriptsize, minimum height=6mm},
  % ---- data variables (circles) ---------------------------------------------
  ppvar/.style={circle, draw=blue!65!black, fill=blue!6, line width=0.6pt,
    inner sep=1pt, minimum size=8mm, font=\footnotesize},
  ppvarA/.style={circle, draw=red!60!black, fill=red!6, line width=0.6pt,
    inner sep=1pt, minimum size=8mm, font=\footnotesize},
  ppvarU/.style={circle, draw=teal!55!black, fill=teal!9, line width=0.6pt,
    inner sep=1pt, minimum size=9mm, font=\scriptsize},
  ppvarY/.style={rounded corners=3pt, draw=blue!60!black, fill=blue!6,
    line width=0.6pt, inner sep=2.5pt, font=\scriptsize},
  % ---- containers -----------------------------------------------------------
  ppplate/.style={rounded corners=6pt, draw=green!45!black, fill=green!5,
    line width=0.7pt},
  ppband/.style={rounded corners=4pt, draw=violet!60!black, fill=none,
    line width=0.6pt, dash pattern=on 1.2pt off 1.6pt},
  % ---- arrows ---------------------------------------------------------------
  ppparr/.style={draw=black!55, line width=0.5pt, dash pattern=on 0.4pt off 1.6pt,
    -{Latex[length=1.5mm,width=1.1mm]}},               % parameter -> step
  ppflow/.style={draw=black!60, line width=0.6pt, -{Latex[length=1.7mm,width=1.2mm]}},
  pparrA/.style={draw=red!55!black, line width=0.5pt, -{Latex[length=1.4mm,width=1mm]}},
  pparrY/.style={draw=blue!55!black, line width=0.5pt, -{Latex[length=1.4mm,width=1mm]}},
  ppdup/.style={draw=teal!55!black, line width=0.6pt, densely dashed,
    -{Latex[length=1.5mm,width=1.1mm]}},
  ppdep/.style={draw=black!40, line width=0.5pt, densely dotted,
    -{Latex[length=1.3mm,width=0.9mm]}},               % gen-step data dependency
  pploop/.style={draw=orange!75!black, line width=0.7pt, dashed,
    dash pattern=on 2.4pt off 1.8pt, -{Latex[length=1.8mm,width=1.3mm]}},
  % ---- a light note attached to a variable ----------------------------------
  ppnote/.style={rounded corners=2.5pt, draw=teal!55!black, fill=teal!4,
    line width=0.5pt, inner sep=4pt, align=left, font=\scriptsize},
]

% =============================================================================
% Background: the K Monte-Carlo replicate plate (L-shaped; notch excludes the
% X column, which is on T rows and only broadcast across K).
% =============================================================================
\draw[ppplate]
  (0.30,8.20) -- (14.05,8.20) -- (14.05,22.45) -- (4.20,22.45)
  -- (4.20,17.90) -- (0.30,17.90) -- cycle;
\node[font=\scriptsize\itshape, color=green!40!black, rotate=-90, anchor=south]
  at (14.35,15.0) {computed on the $T\times K$ Monte-Carlo replicate rows};

% =============================================================================
% 1. Header + the parameter squares that configure the generation steps.
% =============================================================================
\node[pphdr, text width=13.8cm] (hdr) at (7.20,25.85) {%
  \textbf{1.\ Sample dimensions, MLP mechanisms \& noise roots}
  \; {\scriptsize\itshape (once per dataset -- fixed across resamples)}\\[2pt]
  dimensions $d_X,d_U,d_W,d_Z$ with $\min(d_W,d_Z)\!\ge\!d_U$;\quad
  fixed MLPs $f_X,f_U,f_W,f_Z,f_A,f_Y,f_{\mathrm{aug}}$;\quad
  noise $\varepsilon_X,\varepsilon_U,\varepsilon_W,\varepsilon_Z,\varepsilon_A,\varepsilon_Y$};

\node[ppparam] (pX) at (2.30,24.05) {$d_X,\ f_X$};
\node[ppparam] (pU) at (7.00,24.05) {$d_U,\ f_U$};
\node[ppparam] (pW) at (10.60,24.05) {$d_W,\ f_W$};
\node[ppparam] (pZ) at (12.40,24.05) {$d_Z,\ f_Z$};

% the (dimension, mechanism) pairs are the outputs of step 1.
\draw[ppparr] (hdr.south -| pX) -- (pX.north);
\draw[ppparr] (hdr.south -| pU) -- (pU.north);
\draw[ppparr] (hdr.south -| pW) -- (pW.north);
\draw[ppparr] (hdr.south -| pZ) -- (pZ.north);

% =============================================================================
% 2. Generation row.
% =============================================================================
\node[ppgen, text width=2.6cm] (genX)  at (2.30,21.90) {%
  \textbf{2.\ Generate }$X$\\[1pt] $X=f_X(\varepsilon_X)$\ (root)\\[1pt]
  {\scriptsize\itshape on $T$ rows, then broadcast $\times K$}};
\node[ppgen, text width=2.9cm] (genU)  at (7.00,21.90) {%
  \textbf{3.\ Generate }$U$\\[1pt] $U=f_U(X,\varepsilon_U)$};
\node[ppgen, text width=3.4cm] (genWZ) at (11.50,21.90) {%
  \textbf{4.\ Generate }$W,Z$\\[1pt]
  $W=f_W(U,X,5\varepsilon_W)$\\ $Z=f_Z(U,X,5\varepsilon_Z)$};

% parameter -> step (dotted).
\draw[ppparr] (pX.south) -- (genX.north);
\draw[ppparr] (pU.south) -- (genU.north);
\draw[ppparr] (pW.south) -- ([xshift=-9mm]genWZ.north);
\draw[ppparr] (pZ.south) -- ([xshift= 9mm]genWZ.north);

% =============================================================================
% 3. Post-processing BOX: concise caption on TOP, star circles BELOW.
%    Generation arrows terminate on the box's top border; the box maps its four
%    inputs to the post-processed variables X*,U*,W*,Z*.
%    Transparent fill so the plate notch (X* white / U*,W*,Z* green) shows through.
% =============================================================================
\draw[ppband] (0.55,18.20) rectangle (13.75,20.45);
\node[rounded corners=2.5pt, fill=violet!6, draw=violet!45!black, line width=0.4pt,
      inner sep=4pt, align=center, font=\scriptsize, text width=11.6cm,
      text=violet!45!black] (pplab) at (7.15,19.75) {%
  \textbf{5.\ Post-process:} categorise $\cdot$ remove outliers $\cdot$ standardise;
  \ \ $W,Z$ categories \emph{dominate} $U$'s};

\node[ppvar] (Xstar) at (2.30,18.65) {$X^{*}$};
\node[ppvar] (Ustar) at (7.00,18.65) {$U^{*}$};
\node[ppvar] (Wstar) at (10.60,18.65) {$W^{*}$};
\node[ppvar] (Zstar) at (12.40,18.65) {$Z^{*}$};

% generation -> post-process strip -> post-processed variable.
\draw[ppflow] (genX.south)               -- (pplab.north -| genX.south);
\draw[ppflow] (pplab.south -| Xstar)     -- (Xstar.north);
\draw[ppflow] (genU.south)               -- (pplab.north -| genU.south);
\draw[ppflow] (pplab.south -| Ustar)     -- (Ustar.north);
\draw[ppflow] ([xshift=-9mm]genWZ.south) -- (pplab.north -| Wstar);
\draw[ppflow] (pplab.south -| Wstar)     -- (Wstar.north);
\draw[ppflow] ([xshift= 9mm]genWZ.south) -- (pplab.north -| Zstar);
\draw[ppflow] (pplab.south -| Zstar)     -- (Zstar.north);

% =============================================================================
% 4. Mechanism steps: Sample A (left), Sample Y0,Y1 (right); U^aug in between.
% =============================================================================
% 6. U-augmentation (numbered box) -> U^aug -> fed into the A and Y mechanisms.
\node[ppnote, text width=4.9cm] (auglab) at (7.00,16.55) {%
  \textbf{6.\ $U$-augmentation.} $U^{\mathrm{aug}}=f_{\mathrm{aug}}(U^{*})$ --- a
  non-linear expansion of $U$ that increases confounding, fed to $f_A$ and $f_Y$.};
\node[ppvarU] (aug) at (7.00,15.05) {$U^{\mathrm{aug}}$};

\node[ppabox, text width=4.25cm] (A) at (2.95,12.95) {%
  \textbf{7.\ Sample treatment }$A$\\[2pt]
  $\ell=f_A(X,U,Z,U^{\mathrm{aug}},\varepsilon_A)$\\[1pt]
  $A\sim\mathrm{Bern}\!\big(\mathrm{clamp}_\epsilon\,\sigma(a\widetilde\ell+b)\big)$};
\node[ppvarA] (Avar) at (2.95,10.95) {$A$};

\node[ppybox, text width=4.65cm] (Y) at (11.15,12.95) {%
  \textbf{8.\ Sample outcomes }$Y_0,Y_1$\\[2pt]
  $(Y_0,Y_1)=f_Y(X,U,W,U^{\mathrm{aug}},\varepsilon_Y)$\\[1pt]
  \emph{remix:}\ $\gamma$ (CATE variance), $\alpha$ (ATE)};
\node[ppvarY] (Yvar) at (11.15,10.95) {$Y_0,\,Y_1$};

% parent arrows: covariate X* and proxies Z* (->A), W* (->Y); confounder via U^aug.
\draw[pparrA] (Xstar.south) to[out=-90,in=105] (A.north west);
\draw[pparrA] (Zstar.south) to[out=-115,in=45] (A.north east);
\draw[pparrY] (Xstar.south) to[out=-55,in=140] (Y.north west);
\draw[pparrY] (Wstar.south) to[out=-90,in=75]  (Y.north east);

% confounding via U-augmentation: U* -> box 6 -> U^aug -> {A,Y}.
\draw[ppdup] (Ustar.south) -- (auglab.north)
  node[midway,right,font=\scriptsize,inner sep=1.5pt]{$f_{\mathrm{aug}}$};
\draw[ppdup] (auglab.south) -- (aug.north);
\draw[ppdup] (aug.west) to[out=210,in=80]  (A.north);
\draw[ppdup] (aug.east) to[out=-30,in=100] (Y.north);

% step -> output variable.
\draw[ppflow] (A.south) -- (Avar.north);
\draw[ppflow] (Y.south) -- (Yvar.north);

% =============================================================================
% 5/6. Aggregate over K + observed-row draw (inside the plate, at the bottom).
% =============================================================================
\node[ppagg, text width=12.2cm] (agg) at (7.10,9.05) {%
  \textbf{9.\ Aggregate over }$K$\textbf{ \& draw the observed row.}\quad
  $\big(\widehat{\mu}_0(x_i),\widehat{\mu}_1(x_i)\big)=\tfrac1K\sum_{k}\big(Y_{0,i,k},Y_{1,i,k}\big)$;\;
  draw $k_i\!\sim\!\mathcal{U}\{1,\dots,K\}$, set $y_i=a_i\,y_{1}+(1{-}a_i)\,y_{0}$.};

\draw[ppflow] (Avar.south) -- ([xshift=-2.0cm]agg.north);
\draw[ppflow] (Yvar.south) -- ([xshift= 2.0cm]agg.north);

% =============================================================================
% 7. Accept / reject completeness gate (below the plate) + resample loop.
% =============================================================================
\node[ppacc, text width=13.0cm] (acc) at (7.10,7.45) {%
  \textbf{10.\ Accept / reject.}\;
  keep iff $\min(d_W,d_Z)\ge d_U>0$ \emph{and}
  $\mathrm{dom}(W,U)\wedge\mathrm{dom}(Z,U)$ on the realised blocks;
  else resample (give up after \texttt{completeness\_max\_attempts}).};

\draw[ppflow] (agg.south) -- (acc.north);

% resample loop: reject -> back to the generation row.
\draw[pploop] (acc.west) -- ++(-0.25,0) |- (genX.west);
\node[font=\scriptsize\color{orange!62!black}, rotate=90, anchor=south]
  at (-0.02,15.0) {resample the mechanisms};

% =============================================================================
% 8. Output (full width, bottom).
% =============================================================================
\node[ppout, text width=13.0cm] (out) at (7.10,5.95) {%
  \textbf{11.\ Output.}\;
  context $\{X,\,U_k,\,W_k,\,Z_k,\,A_k\}$ with observed $Y_k(A_k)$;\quad
  regression targets $\mu_0,\mu_1$ (the $K$-averaged potential outcomes from step 9).};

\draw[ppflow] (acc.south) -- (out.north);

\end{tikzpicture}
